\section{Pembuktian dengan Kontraposisi}

Tidak jarang pembuktian langsung menghadapi hambatan teknis --- misalnya ketika manipulasi aljabar dari hipotesis $p$ menuju kesimpulan $q$ menjadi rumit atau tidak alami. Dalam situasi tersebut, pembuktian dengan \logic{kontraposisi} menjadi alternatif yang sangat efektif. Metode ini dilandasi oleh fakta bahwa implikasi $p \rightarrow q$ secara logis ekuivalen dengan kontraposisinya $\neg q \rightarrow \neg p$ \cite{rosen2019,forallx,libretexts_proofs}.

\subsection{Landasan Logis}

Ekuivalensi $p \rightarrow q \equiv \neg q \rightarrow \neg p$ dapat diverifikasi melalui tabel kebenaran:

\begin{center}
\renewcommand{\arraystretch}{1.2}
\begin{tabular}{|c|c||c|c|c|c||c|}
\hline
$p$ & $q$ & $\neg p$ & $\neg q$ & $p \rightarrow q$ & $\neg q \rightarrow \neg p$ & Sama? \\ \hline
T & T & F & F & T & T & Ya \\ \hline
T & F & F & T & F & F & Ya \\ \hline
F & T & T & F & T & T & Ya \\ \hline
F & F & T & T & T & T & Ya \\ \hline
\end{tabular}
\end{center}

Karena kedua kolom selalu menghasilkan nilai yang identik, keduanya merupakan ekspresi yang ekuivalen secara logis. Membuktikan kontraposisi $\neg q \rightarrow \neg p$ secara otomatis membuktikan implikasi asli $p \rightarrow q$.

\subsection{Struktur Bukti Kontraposisi}

Untuk membuktikan $p \rightarrow q$ melalui kontraposisi:

\begin{enumerate}
    \item \textbf{Asumsikan} $\neg q$ (negasi dari kesimpulan) bernilai benar.
    \item \textbf{Gunakan} penalaran deduktif berdasarkan definisi, aksioma, dan teorema yang berlaku.
    \item \textbf{Tunjukkan} bahwa $\neg p$ (negasi dari hipotesis) bernilai benar.
\end{enumerate}

\subsection{Kapan Menggunakan Kontraposisi?}

Kontraposisi sangat berguna dalam situasi-situasi berikut:
\begin{itemize}
    \item Kesimpulan $q$ mengandung negasi atau kondisi ``tidak'' yang sulit diekspresikan secara positif.
    \item Hipotesis $p$ bersifat kompleks sehingga sulit untuk memulai manipulasi aljabar darinya.
    \item Negasi dari kesimpulan ($\neg q$) memberikan informasi yang lebih mudah diolah dibandingkan hipotesis aslinya.
\end{itemize}

\subsection{Contoh: Paritas \texorpdfstring{$3n+2$}{3n+2} dan \texorpdfstring{$n$}{n}}

\begin{contoh}
\textbf{Teorema}: Jika $3n+2$ adalah ganjil, maka $n$ adalah ganjil.

\textbf{Mengapa bukti langsung sulit?} Dari $3n+2 = 2k+1$ (ganjil), didapat $n = \frac{2k-1}{3}$. Pembagian oleh 3 menyulitkan pembuktian bahwa hasilnya ganjil.

\textbf{Bukti (Kontraposisi):} Buktikan kontraposisi: ``Jika $n$ genap, maka $3n+2$ genap.''

\begin{enumerate}
    \item Asumsikan $n$ genap ($\neg q$). Maka $n = 2k$ untuk suatu bilangan bulat $k$.
    \item Hitung $3n + 2$:
    \[ 3n + 2 = 3(2k) + 2 = 6k + 2 = 2(3k + 1) \]
    \item Misalkan $m = 3k + 1$. Karena $k$ bilangan bulat, maka $m$ bilangan bulat.
    \item Sehingga $3n + 2 = 2m$, yang berbentuk bilangan genap ($\neg p$).
    \item Kontraposisi terbukti, maka teorema asli juga terbukti. $\blacksquare$
\end{enumerate}
\end{contoh}

\subsection{Contoh: Kuadrat Genap Implies Bilangan Genap}

\begin{contoh}
\textbf{Teorema}: Jika $n^2$ genap, maka $n$ genap.

\textbf{Bukti (Kontraposisi):} Buktikan: ``Jika $n$ ganjil, maka $n^2$ ganjil.''

\begin{enumerate}
    \item Asumsikan $n$ ganjil. Maka $n = 2k + 1$ untuk suatu bilangan bulat $k$.
    \item Hitung $n^2$:
    \[ n^2 = (2k + 1)^2 = 4k^2 + 4k + 1 = 2(2k^2 + 2k) + 1 \]
    \item Misalkan $m = 2k^2 + 2k$, yang merupakan bilangan bulat.
    \item Sehingga $n^2 = 2m + 1$, yang berbentuk bilangan ganjil.
    \item Kontraposisi terbukti: $n$ ganjil $\Rightarrow$ $n^2$ ganjil. Maka kontraposisinya, ``$n^2$ genap $\Rightarrow$ $n$ genap,'' juga terbukti. $\blacksquare$
\end{enumerate}
\end{contoh}

\begin{catatan}
Perhatikan bahwa contoh kedua (``$n^2$ genap $\Rightarrow$ $n$ genap'') sangat sulit dibuktikan secara langsung karena dari $n^2 = 2k$ kita memperoleh $n = \sqrt{2k}$, dan tidak ada cara langsung untuk menyimpulkan bahwa $\sqrt{2k}$ genap. Kontraposisi menghindari akar kuadrat dan membuat bukti mengalir dengan natural.
\end{catatan}
